\documentclass[aps,pre,twocolumn,superscriptaddress,floatfix]{revtex4-2}
\usepackage{amsmath,amssymb}
\usepackage{graphicx}
\usepackage{hyperref}

\begin{document}

\title{Paper~3: Temporal Dynamics of Impedance Networks\\
--- Impedance Debt, Sleep, Lifecycle, and Aging}

\author{Hsi-Yu Huang}
\email{alice.gamma.net@gmail.com}
\affiliation{Independent Researcher, Taiwan}

\date{\today}

\begin{abstract}
Papers~1 and~2 of this series established the spatial
architecture of impedance-minimising networks
(relay-node emergence, fractal scaling, Murray's Law,
dual neural--vascular organ health
$H = (1-\Gamma_n^2)(1-\Gamma_v^2)$).
Here we develop the full \emph{temporal} theory in two
complementary parts.

\textbf{Part~A} introduces \emph{impedance debt}
$D_Z = \int |\Gamma|^2 P_{\mathrm{in}}\,dt$:
the time-integrated mismatch energy that accumulates
during wakefulness.
We show that sleep is the minimum-action mechanism
for discharging $D_Z$---a thermodynamic necessity,
not a behavioural choice.
Decomposing $D_Z$ into microscopic (synaptic drift)
and thermal (waste heat) components explains why
jellyfish sleep without a brain
($D_{\mathrm{thermal}}$ removed by ocean)
and why brains evolved convergently upon the water-to-land
transition ($D_{\mathrm{thermal}}$ must be managed internally).
A dual-field extension ($Z$-field + material field $\rho$)
unifies Wolff's law (bone), Davis's law (muscle),
and impedance-remodeling (neurons) as parameter
instances of a single impedance-matching PDE,
with embryogenesis recapitulating the thermodynamic
causal chain (heart $\to$ neural tube $\to$ brain).

\textbf{Part~B} derives the \emph{Lifecycle Equation},
\begin{equation*}
  \frac{d(\Sigma\Gamma^2)}{dt}
    = \underbrace{-\eta\,\Sigma\Gamma^2}_{\text{learning}}
    + \underbrace{\gamma\,\Gamma_{\mathrm{env}}(t)}_{\text{novelty}}
    + \underbrace{\delta\,D(t)}_{\text{aging}}\,,
\end{equation*}
governing the complete cognitive trajectory from birth to
senescence as a three-force competition on the dual
neural--vascular network.
The resulting \emph{bathtub curve} of $\Sigma\Gamma^2$
reproduces the known developmental trajectory without
any fitted function.
Emotion, reward, stress, curiosity, and boredom emerge as
\emph{readouts} of $\Gamma$ dynamics at different time scales:
dopamine encodes $-\Delta\Gamma^2$, cortisol encodes chronic
$\Gamma^2$ elevation, curiosity encodes the impedance gradient,
and boredom encodes stagnation at non-zero mismatch.
Aging is derived as Coffin--Manson fatigue accelerated by
Arrhenius stress kinetics, with
atherosclerosis as vascular fatigue and the cascade
$\Gamma_v\!\uparrow \to \rho\!\downarrow \to \Gamma_n\!\uparrow$
explaining why cardiovascular disease precedes
neurodegeneration.
Eighteen computational experiments verify these predictions.
\end{abstract}

\maketitle


% %%%%%%%%%%%%%%%%%%%%%%%%%%%%%%%%%%%%%%%%%%%%%%%%%%%%%%%%%%%%
%  PART A — IMPEDANCE DEBT AND SLEEP
% %%%%%%%%%%%%%%%%%%%%%%%%%%%%%%%%%%%%%%%%%%%%%%%%%%%%%%%%%%%%
\part*{Part A: Impedance Debt and Sleep}
\label{part:A}


% ============================================================
%  SECTION 1 — INTRODUCTION
% ============================================================
\section{Introduction}
\label{sec:intro}

Paper~1~\cite{Paper1} showed that the Minimum
Reflection Action $\mathcal{A}[\Gamma] = \int |\Gamma|^2\,dt$
on an impedance-heterogeneous network forces hierarchical
relay-node emergence, with dimensional cut cost $A_{\mathrm{cut}}$
irreducible by any local rule; fractal scaling recovers
Kleiber's $B \sim M^{3/4}$ and Murray's branching exponent.
Paper~2~\cite{Paper2} extended the framework to vascular
networks and established the dual-network organ model
$H = (1-\Gamma_n^2)(1-\Gamma_v^2)$.

These results are \emph{spatial}: they describe the
equilibrium architecture of impedance networks.
But biological networks operate in time.
Action potentials flow, ion channels drift,
metabolic heat accumulates.
What is the cost of running the network?

This paper extends the Minimum Reflection Action to the
time domain.
Part~A defines impedance debt $D_Z$---the integral of
$|\Gamma|^2$ weighted by incident power---and shows that
sleep is its minimum-action discharge mechanism.
The decomposition of $D_Z$ into microscopic and thermal
components explains cnidarian sleep, convergent brain
evolution, and embryogenesis as a single thermodynamic
causal chain.
Part~B derives the Lifecycle Equation---the first-order ODE
governing $\Sigma\Gamma^2$ from birth to death---and shows
that cognition, emotion, and aging are three forces competing
on the same impedance network.


% ============================================================
%  SECTION 2 — PHYSICAL FRAMEWORK
% ============================================================
\section{Physical Framework}
\label{sec:framework}

\subsection{Impedance mismatch and the three constraints}

At every signalling interface, the reflection coefficient is
\begin{equation}
  \Gamma = \frac{Z_{\mathrm{load}} - Z_{\mathrm{source}}}
               {Z_{\mathrm{load}} + Z_{\mathrm{source}}}\,.
  \label{eq:gamma}
\end{equation}
The three inviolable constraints are:
\begin{itemize}
  \item \textbf{C1 (Energy Conservation)}:
    $\Gamma^2 + T = 1$ at every interface, every tick.
  \item \textbf{C2 (Impedance Remodeling)}:
    $\Delta Z = -\eta\,\Gamma\,x_{\mathrm{in}}\,x_{\mathrm{out}}$.
  \item \textbf{C3 (Impedance-Tagged Transport)}:
    All inter-module signals are \texttt{ElectricalSignal}
    objects carrying $Z$-metadata.
\end{itemize}


\subsection{Neurons as transmission lines}
\label{sec:neuron-TL}

A myelinated axon has characteristic impedance
\begin{equation}
  Z_0 = \sqrt{\frac{R + j\omega L}{G + j\omega C}}\,,
  \label{eq:Z0}
\end{equation}
where $R$, $L$, $G$, $C$ are the distributed per-unit-length
resistance, inductance, conductance, and capacitance
of the neural membrane.
At the synapse, $Z_0$ meets the postsynaptic input impedance
$Z_{\mathrm{post}}$, producing a reflection coefficient
$\Gamma = (Z_{\mathrm{post}} - Z_0)/(Z_{\mathrm{post}} + Z_0)$.
Perfect transmission ($\Gamma = 0$) requires impedance matching;
any mismatch wastes energy as reflected power.

Over time, ion-channel conductances drift (fatigue,
temperature fluctuation, molecular turnover), shifting both
$Z_0$ and $Z_{\mathrm{post}}$ away from their initial values.
This drift accumulates as \emph{impedance debt}.


\subsection{Definition of impedance debt}
\label{sec:DZ-def}

Define the impedance debt at time $t$ as
\begin{equation}
  \boxed{\;
  D_Z(t) = \int_0^{t} \frac{1}{N}
    \sum_{i=1}^{N} |\Gamma_i(t')|^2\,
    P_{\mathrm{in},i}(t')\,dt'
  \;}
  \label{eq:DZ}
\end{equation}
where $N$ is the number of signalling interfaces and
$P_{\mathrm{in},i}$ is the incident power at interface~$i$.
$D_Z$ has units of energy and is monotonically
non-decreasing during wakefulness
(active signalling guarantees $P_{\mathrm{in}} > 0$).

When the impedance landscape changes slowly relative to
firing rate, $\langle|\Gamma|^2\rangle$ is approximately
constant over a wake epoch.
Then $D_Z$ grows linearly:
\begin{equation}
  D_Z(t) \approx \langle|\Gamma|^2\rangle\,
    \bar{P}_{\mathrm{in}}\,t\,.
  \label{eq:DZ-linear}
\end{equation}


\subsection{Adenosine as a $D_Z$ readout}
\label{sec:adenosine}

Adenosine accumulates extracellularly in proportion to
metabolic demand---specifically, to ATP hydrolysis rate,
which is proportional to $P_{\mathrm{dissipated}} =
\Gamma^2 \cdot P_{\mathrm{in}}$.
Thus extracellular adenosine concentration
$[\mathrm{Ade}](t) \propto D_Z(t)$.
This identifies adenosine as a \emph{local impedance-debt
meter}---consistent with its established role as the
primary molecular correlate of sleep pressure
(Borb\'{e}ly's Process~S~\cite{Borbely1982}).


\subsection{$D_Z$ decomposition: microscopic and thermal}
\label{sec:DZ-decomp}

The total impedance debt decomposes into two physically
distinct components:
\begin{equation}
  D_Z = D_{\mathrm{micro}} + D_{\mathrm{thermal}}\,.
  \label{eq:DZ-decomp}
\end{equation}
$D_{\mathrm{micro}}$ (synaptic impedance drift) is the
cumulative mismatch from ion-channel stochastic drift,
receptor desensitisation, and molecular turnover.
$D_{\mathrm{thermal}}$ (waste heat) is the $\Gamma^2$-derived
dissipation that raises local tissue temperature.
This decomposition has evolutionary consequences:
organisms immersed in high-$\kappa$ media (ocean) can
passively dissipate $D_{\mathrm{thermal}}$ but must still
manage $D_{\mathrm{micro}}$ via sleep-like quiescence.


\subsection{Dual-network impedance debt}
\label{sec:dual-network-debt}

Paper~2~\cite{Paper2} established that organ health depends
on two impedance networks.
Impedance debt likewise splits:
\begin{equation}
  D_Z = D_{Z,n} + D_{Z,v}\,,
  \label{eq:DZ-dual}
\end{equation}
where $D_{Z,n}$ is neural and $D_{Z,v}$ is vascular debt.
The discharge time constants differ by orders of magnitude:
neural debt discharges on the timescale of a night's sleep
($\eta_n \sim 0.08$);
vascular debt discharges on weeks to months
($\eta_v \sim 0.001$).
This asymmetry means nightly sleep primarily services
$D_{Z,n}$, while $D_{Z,v}$ accumulates chronically
unless offset by vascular remodeling.
The coupling from Paper~2,
$I_{\mathrm{blood}} = (1-\Gamma_v^2)\,Q_0$,
means rising vascular debt ($\Gamma_v\!\uparrow$) starves
neural repair of material ($\rho\!\downarrow$),
amplifying neural debt in a positive-feedback loop.


\subsection{Sleep as impedance-debt discharge}
\label{sec:sleep}

During NREM sleep, the brain activates a recalibration
mode: reduced sensory input ($P_{\mathrm{in}} \to 0$)
halts $D_Z$ accumulation, while slow oscillations drive
$Z$-map restoration via C2 impedance remodeling.
The discharge rate follows:
\begin{equation}
  D_Z(t) = D_Z^{\mathrm{peak}}\,
    e^{-\beta_{\mathrm{recal}}\,(t - t_{\mathrm{sleep}})}\,,
  \label{eq:DZ-decay}
\end{equation}
where $\beta_{\mathrm{recal}}$ is the N3 recalibration rate.
The organism sleeps until $D_Z < D_Z^*$ (threshold).

The maximum tolerable wake duration before $D_Z$ exceeds
a critical threshold $D_Z^*$ is:
\begin{equation}
  \boxed{\;
  T_{\mathrm{wake}}^{\max}
    = \frac{D_Z^*}
           {\langle|\Gamma|^2\rangle\,\bar{P}_{\mathrm{in}}}
  \;}
  \label{eq:Twake}
\end{equation}
This explains why infants ($\langle\Gamma^2\rangle \approx 0.12$)
sleep $\sim$16\,h/day while adults
($\langle\Gamma^2\rangle \approx 0.04$) need only $\sim$8\,h:
lower average mismatch means slower debt accumulation.


\subsection{NREM/REM partition}
\label{sec:NREM-REM}

NREM sleep discharges $D_Z$ (impedance recalibration).
REM sleep reorganises the $Z$-map topology
(adding/pruning connections).
The required REM fraction scales with topological complexity:
\begin{equation}
  T_{\mathrm{REM}} \propto N_{\mathrm{topo}}\,,
  \label{eq:REM-fraction}
\end{equation}
where $N_{\mathrm{topo}}$ is the number of active nodes
in the impedance network.
This explains the high REM fraction in neonates
($\sim$50\%, rapidly growing $N$) versus adults
($\sim$25\%, stable $N$).


% ============================================================
%  SECTION 3 — DUAL-FIELD EXTENSION
% ============================================================
\section{Dual-Field Extension: Morphogenesis}
\label{sec:dual-field}

When impedance varies in space (tissue structure)
and is coupled to a material-supply field,
the framework becomes a reaction--diffusion system:
\begin{align}
  \frac{\partial Z}{\partial t}
    &= D_Z\,\nabla^2 Z
    \underbrace{-\eta\,f(\rho)\,\Gamma\,Z}_{\text{C2 remodeling}}
    \underbrace{-\,E(\Gamma^2)\,(-\Gamma)\,\rho}_{\text{enzyme}}
    - \lambda\,Z\,,
  \label{eq:Z-field} \\[4pt]
  \frac{\partial\rho}{\partial t}
    &= D_\rho\,\nabla^2\rho
    + I_{\mathrm{blood}}
    - c(\Gamma^2)\,\rho\,,
  \label{eq:rho-field}
\end{align}
where:
\begin{itemize}
  \item $D_Z$, $D_\rho$ are diffusion coefficients,
  \item $f(\rho) = \rho/(\rho + K_\rho)$: material-limited
    repair (Michaelis--Menten form),
  \item $E(\Gamma^2)$: Hill enzyme kinetics,
  \item $I_{\mathrm{blood}} = (1-\Gamma_v^2)\,Q_0$: vascular
    supply from Paper~2,
  \item $\lambda$: entropic degradation rate.
\end{itemize}

The factor $(-\Gamma)$ in the enzyme term ensures
bidirectional remodeling: building when $Z < Z_0$
($\Gamma < 0$, osteoblast-like) and degrading when
$Z > Z_0$ ($\Gamma > 0$, osteoclast-like)---the
mathematical expression of Wolff's law~\cite{Wolff1892}.
Under appropriate parameters, the same PDE reduces to
Davis's law (muscle) and impedance remodeling (neurons).

When the diffusion ratio $D_\rho / D_Z$ exceeds a
tissue-specific threshold, the uniform steady state
becomes unstable via Turing instability~\cite{Turing1952},
producing spatial patterns in the $Z$-field
(bone trabecular architecture, hepatic lobule periodicity).

The material-bottleneck limit ($\rho \to 0$) yields:
\begin{equation}
  \rho \to 0 \;\Longrightarrow\;
  \frac{\partial Z}{\partial t} \to -\lambda\,Z\,.
  \label{eq:ischemia-limit}
\end{equation}
Only entropic degradation survives---no repair is possible
without material supply.


% ============================================================
%  SECTION 4 — EVIDENCE I: CNIDARIAN SLEEP
% ============================================================
\section{Evidence I: Cnidarian Sleep Without a Brain}
\label{sec:cnidaria}

\textit{Cassiopea} satisfies all sleep criteria
(reduced activity, increased arousal threshold,
homeostatic rebound~\cite{Nath2017}).
Appelbaum \emph{et al.}~\cite{Appelbaum2026} showed
that neuronal DNA double-strand breaks accumulate linearly
with wake duration and are repaired during sleep.

In the impedance-debt framework, DNA damage is a
\emph{downstream marker} of $D_Z$:
the causal chain is
$\Gamma^2\!\uparrow \to T\!\downarrow \to$
mitochondrial compensation $\to$ ROS $\to$ DNA damage.
Jellyfish, immersed in an ocean with volumetric heat
capacity $\sim$3\,500$\times$ that of air,
passively dissipate $D_{\mathrm{thermal}}$.
Yet $D_{\mathrm{micro}}$ (ion-channel drift) still
accumulates, requiring sleep-like quiescence for discharge.
Hence: \textbf{no brain is needed, but sleep is still required.}


% ============================================================
%  SECTION 5 — EVIDENCE II: CONVERGENT BRAIN EVOLUTION
% ============================================================
\section{Evidence II: Convergent Evolution of Brains}
\label{sec:brains}

Complex nervous systems evolved independently at least
nine times across Metazoa~\cite{Moroz2012}.
We argue this reflects a common thermodynamic pressure
traceable to one event: saturation of the filter-feeding
niche in the late Ediacaran.

\subsection{The Cambrian phase transition}

For two billion years, filter-feeding was optimal:
the impedance to nutrient sources was low
($Z_{\mathrm{food}} \approx 1/(\rho_{\mathrm{micro}}
D_{\mathrm{diff}} A_{\mathrm{oral}})$).
Population overshoot drove $\rho_{\mathrm{micro}}$
down, sending $Z_{\mathrm{food}} \to \infty$ and
$\Gamma_{\mathrm{feeding}}^2 \to 1$.
The minimum-action escape was predation---a new channel
with lower $\Gamma^2$ but requiring eyes, appendages,
and coordinated impedance management (i.e., a nervous system).
The resulting arms race drove the Cambrian explosion:
each cycle of attack and defense increased both channel
count and irreducible cut cost $A_{\mathrm{cut}}$ (Paper~1).

\subsection{Water-to-land transition}

When oceanic niches saturated, displaced lineages colonised
land---exchanging seawater ($\kappa \ge 1\;\mathrm{W\,m^{-1}K^{-1}}$)
for air ($\kappa \le 0.04$).
$D_{\mathrm{thermal}}$ could no longer be passively removed,
creating selective pressure for:
(i)~a convective heat-transport network (vasculature---Paper~2),
(ii)~a centralised controller (brain),
(iii)~structured sleep architecture (NREM/REM).
Cephalopods remained marine but evolved high metabolic
rates that \emph{exceeded} oceanic heat-sink capacity,
solving the problem via distributed brain architecture
(nine ganglia)---confirmed by Experiment~A7 at
$\kappa_{\mathrm{cross}} \approx 0.30$.

The tidal zone acts as a quarter-wave impedance matcher:
$Z_{\mathrm{match}} = \sqrt{Z_w Z_l}$, with a logarithmic
$\kappa$ gradient reducing the transition action
$\mathcal{A}[\Gamma]$ to 6.2\% of a direct jump
(Experiment~A6).
This explains why the water-to-land transition occurred
independently at least five times---it is the
minimum-action evolutionary path.


% ============================================================
%  SECTION 6 — EVIDENCE III: EMBRYOGENESIS
% ============================================================
\section{Evidence III: Embryogenesis}
\label{sec:embryo}

Human embryonic development follows the thermodynamic
causal chain:
\begin{enumerate}
  \item \textbf{Heart (day 22--23)}: establishes the vascular
    impedance network ($I_{\mathrm{blood}} = (1-\Gamma_v^2)Q_0$).
  \item \textbf{Neural tube (week 3--4)}: once vascular supply
    provides $\rho > 0$, high-metabolic neural tissue can form.
  \item \textbf{Brain vesicles (week 5--6)}: centralised
    impedance-debt manager for both networks.
\end{enumerate}
Each step is a thermodynamic prerequisite for the next.
The fetus floats in amniotic fluid at $\sim$37.6$^\circ$C,
with $\sim$85\% of metabolic heat removed via the placenta
---thermodynamically a jellyfish whose heat sink is the mother.
\textbf{Birth is the water-to-land transition replayed
in ontogeny.}


% %%%%%%%%%%%%%%%%%%%%%%%%%%%%%%%%%%%%%%%%%%%%%%%%%%%%%%%%%%%%
%  PART B — LIFECYCLE EQUATION AND AGING
% %%%%%%%%%%%%%%%%%%%%%%%%%%%%%%%%%%%%%%%%%%%%%%%%%%%%%%%%%%%%
\part*{Part B: Lifecycle Equation and Aging}
\label{part:B}


% ============================================================
%  SECTION 7 — THE LIFECYCLE EQUATION
% ============================================================
\section{The Lifecycle Equation}
\label{sec:lifecycle}

\subsection{Total mismatch energy}

Define the population-averaged squared reflection coefficient:
\begin{equation}
  \Sigma\Gamma^2(t) \equiv
    \frac{1}{N} \sum_{i=1}^{N} \Gamma_i^2(t)\,,
  \label{eq:sigma-gamma}
\end{equation}
with dual-network decomposition
$\Sigma\Gamma^2 = \Sigma\Gamma_n^2 + w_v\,\Sigma\Gamma_v^2$,
where $w_v = \alpha_{v\to n}$ is the vascular-to-neural
coupling weight from Paper~2.

\subsection{Three forces}

Three physically distinct mechanisms alter $\Sigma\Gamma^2$:

\emph{(i)~Learning (MRP descent).}
C2 impedance remodeling adapts $Z_{\mathrm{load}}$ toward
$Z_{\mathrm{source}}$:
\begin{equation}
  \left.\frac{d(\Sigma\Gamma^2)}{dt}\right|_{\mathrm{learn}}
    = -\eta(t)\,\Sigma\Gamma^2\,.
  \label{eq:learning-term}
\end{equation}

\emph{(ii)~Novelty injection.}
New stimuli introduce mismatches at untuned interfaces:
\begin{equation}
  \left.\frac{d(\Sigma\Gamma^2)}{dt}\right|_{\mathrm{novel}}
    = \gamma(t)\,\Gamma_{\mathrm{env}}(t)\,.
  \label{eq:novelty-term}
\end{equation}

\emph{(iii)~Aging (Coffin--Manson fatigue).}
Irreversible mechanical strain from Lorentz pinch
($\varepsilon_{\mathrm{pinch}} = S|I|^\alpha$)
drifts impedance permanently, with stress-accelerated
kinetics:
\begin{equation}
  \delta_{\mathrm{eff}}(t) = \delta_0\,
    \exp\!\bigl(E_a\,\sigma_{\mathrm{stress}}(t)\bigr)\,.
  \label{eq:arrhenius-aging}
\end{equation}

Combining all three:
\begin{equation}
  \boxed{\;
  \frac{d(\Sigma\Gamma^2)}{dt}
    = -\eta\,\Sigma\Gamma^2
    + \gamma\,\Gamma_{\mathrm{env}}(t)
    + \delta\,D(t)
  \;}
  \label{eq:lifecycle}
\end{equation}
This is a non-autonomous first-order ODE whose coefficients
are functions of developmental stage, environmental
statistics, and cumulative damage.


\subsection{The Equilibrium Equation}

Impedance mismatch dissipates heat.
Thermal balance gives the steady-state temperature:
\begin{equation}
  T_{\mathrm{steady}} = T_{\mathrm{env}}
    + \frac{\alpha}{\beta}\,\Sigma\Gamma^2\,.
  \label{eq:equilibrium}
\end{equation}
``Adulthood'' has a physical definition:
the thermal steady state where $d(\Sigma\Gamma^2)/dt \approx 0$
---the point at which the world can no longer burn you.


% ============================================================
%  SECTION 8 — EMOTION AS IMPEDANCE DYNAMICS
% ============================================================
\section{Emotion as Impedance Dynamics}
\label{sec:emotion}

Psychological states map to features of the $\Gamma^2$
time series:

\subsection{Dopamine: the $-\Delta\Gamma^2$ signal}

The reward prediction error~\cite{Schultz1997} in
the impedance framework is:
\begin{equation}
  \delta_{\mathrm{RPE}} = -\bigl(\Sigma\Gamma^2_{t+1}
    - \mathbb{E}[\Sigma\Gamma^2_{t+1}]\bigr)\,.
  \label{eq:rpe}
\end{equation}
Dopamine fires when mismatch decreases \emph{more than
expected}---the impedance network matched faster
than predicted.

\subsection{Cortisol: the slow $\Gamma^2$ integrator}

Stress is the time-averaged mismatch:
\begin{equation}
  \sigma_{\mathrm{stress}}(t) =
    \frac{1}{\tau_s} \int_{t-\tau_s}^{t}
    \Sigma\Gamma^2(t')\,dt'\,,
  \label{eq:cortisol}
\end{equation}
where $\tau_s \sim$ hours--days.
Acute challenges are stimulating; prolonged elevation
activates the cortisol response.

\subsection{The Yerkes--Dodson curve}

Arousal increases both $\eta$ (learning) and $\gamma$
(novelty intake).
The optimum occurs at:
\begin{equation}
  \eta^*\,\Sigma\Gamma^2
    = \gamma^*\,\Gamma_{\mathrm{env}}
  \quad\Longrightarrow\quad
  \text{arousal}^* \propto
    \frac{\Sigma\Gamma^2}{\Gamma_{\mathrm{env}}}\,,
  \label{eq:yerkes-dodson}
\end{equation}
predicting that optimal arousal depends on task difficulty
and current skill level---the Yerkes--Dodson
observation~\cite{YerkesDodson1908}.

\subsection{Curiosity and boredom}

Curiosity is the impedance gradient:
$\mathcal{C} = -\partial(\Sigma\Gamma^2)/\partial(\text{exploration})$,
gated by safety ($\Gamma^2 < \Gamma^2_{\mathrm{threat}}$).
Boredom arises when $|d\Gamma^2/dt| < \epsilon$ at
non-zero mismatch---a thermodynamic imperative to seek
new gradients for action reduction.

\subsection{Metacognition}

The thinking impedance
$\Gamma_{\mathrm{think}} = \sum_i w_i\Gamma_i / \sum_i w_i$
defines dual-process switching~\cite{Kahneman2011}:
System~1 (fast, automatic) when
$\Gamma_{\mathrm{think}} < \Gamma_{\mathrm{S2}}$;
System~2 (slow, deliberate) when
$\Gamma_{\mathrm{think}} > \Gamma_{\mathrm{S2}}$.


% ============================================================
%  SECTION 9 — THE BATHTUB CURVE
% ============================================================
\section{The Bathtub Curve of Cognition}
\label{sec:bathtub}

The Lifecycle Equation predicts six developmental phases:
\begin{enumerate}
  \item \textbf{Birth}: $\Sigma\Gamma^2 \approx 0.8$--$1.0$.
    Random initialisation; Novelty $\gg$ Learning $\gg$ Aging.
  \item \textbf{Infancy}: Learning dominates.
    $\eta$ maximal (critical period); $\Sigma\Gamma^2$ drops rapidly.
  \item \textbf{Childhood}: Learning $\approx$ Novelty.
    Intense consolidation; REM fraction high.
  \item \textbf{Adulthood}: $d(\Sigma\Gamma^2)/dt \approx 0$.
    Equilibrium Equation gives steady-state temperature.
  \item \textbf{Decline}: Aging $>$ Learning.
    Net mismatch increases despite continued learning.
  \item \textbf{Senescence}: Aging $\gg$ Learning.
    Cascading failure with positive feedback.
\end{enumerate}
This trajectory is formally identical to the engineering
reliability \emph{bathtub curve}~\cite{Klutke2003}---not
analogy but derivation from the same physics.


% ============================================================
%  SECTION 10 — STRESS, FATIGUE, AND AGING
% ============================================================
\section{Stress, Fatigue, and the Physics of Aging}
\label{sec:aging}

\subsection{Coffin--Manson fatigue}

Action potentials generate Lorentz forces; cumulative
strain follows the Coffin--Manson law~\cite{Coffin1954}:
$N_f = C\,(\Delta\varepsilon_p)^{-\beta}$.
``Failure'' is irreversible impedance drift.

\subsection{Vascular Coffin--Manson: atherosclerosis}

$\sim$3 billion cardiac cycles impose pulsatile strain.
Cumulative vascular damage $D_v(t)$ drifts $Z_v$ from
Murray's-Law optimum (Paper~2), with plaque as the
macroscopic manifestation:
a 50\% stenosis raises $Z_v$ by $16\times$ via
$r^{-4}$ Poiseuille scaling.

\subsection{The vascular--neural aging cascade}

\begin{equation}
  D_v\!\uparrow \;\to\;
    \Gamma_v\!\uparrow \;\to\;
    \rho\!\downarrow \;\to\;
    \eta_{n,\mathrm{eff}}\!\downarrow \;\to\;
    \Gamma_n\!\uparrow\,.
  \label{eq:aging-cascade}
\end{equation}
This explains:
(i)~cardiovascular disease precedes dementia by 10--20 years,
(ii)~blood-pressure control in midlife reduces dementia risk
(SPRINT-MIND~\cite{sprint2019}),
(iii)~metabolic syndrome as dual-network failure.

\subsection{Elastic versus plastic deformation}

Each current pulse produces $\varepsilon = \varepsilon_{\mathrm{elastic}}
+ \varepsilon_{\mathrm{plastic}}$.
Sleep repairs the elastic component (daily $D_Z$ discharge);
the plastic component is \emph{irreversible}---it accumulates
over a lifetime.
Sleep restores daily performance but cannot prevent
long-term aging~\cite{Walker2017}.

\subsection{Death as positive-feedback collapse}

When $\Sigma\Gamma^2$ exceeds a critical threshold:
\begin{equation}
  \frac{d^2(\Sigma\Gamma^2)}{dt^2} > 0
  \quad\Longrightarrow\quad
  \text{runaway failure}\,.
  \label{eq:runaway}
\end{equation}
Compensatory current increases through remaining channels
accelerate fatigue in a positive-feedback loop,
amplified by the dual-network cascade.


% ============================================================
%  SECTION 11 — SLEEP IN THE LIFECYCLE CONTEXT
% ============================================================
\section{Sleep Across the Lifecycle}
\label{sec:sleep-lifecycle}

\emph{Infant sleep:}
$\Sigma\Gamma^2 \approx 0.8$; daily $D_Z$ is enormous.
14--17 hours of sleep required.
REM $\sim$50\% because $N_{\mathrm{topo}}$ is rapidly growing.

\emph{Adult sleep:}
$\Sigma\Gamma^2$ near minimum; 7--9 hours suffice.

\emph{Elderly sleep:}
Irreversible (plastic) drift is undischargeable by sleep.
Rising vascular $\Gamma_v$ (atherosclerosis) reduces
material supply $\rho$, slowing neural repair.
Sleep is doubly impaired: plastic baseline rises
\emph{and} vascular-mediated repair
declines~\cite{Mander2017}.


% ============================================================
%  SECTION 12 — INFANT MECHANICS AND SOCIAL CALIBRATION
% ============================================================
\section{Developmental Design Principles}
\label{sec:developmental}

\subsection{Infant soft tissue as impedance buffer}

Neonatal soft tissues (cartilage, fat, elastic skin) form
\emph{mechanical quarter-wave transformers}:
$Z_{\text{infant}} \ll Z_{\text{adult}}$,
attenuating impact signals before they reach nociceptors.
With $\Sigma\Gamma^2 \approx 0.8$ at birth,
unattenuated pain would overwhelm C2 calibration.

\subsection{Social referencing as external $\Gamma$ calibration}

When an infant falls, it looks at the caregiver's face
before deciding to cry~\cite{Sorce1985}.
The caregiver's emotional signal provides
$\Gamma_{\mathrm{social}}$ superimposed on the
mechanical signal:
$\Gamma_{\mathrm{perceived}} =
\Gamma_{\mathrm{mechanical}} + \eta_{\mathrm{social}}\,
\Gamma_{\mathrm{social}}$,
where $\eta_{\mathrm{social}} \gg 1$ in infancy
and $\to 0$ in adulthood (self-calibrated system).
Each fall-and-look episode is a C2 calibration step:
$\Delta Z_{\text{pain}} = -\eta\,\Gamma_{\text{social}}\,
x_{\mathrm{in}}\,x_{\mathrm{out}}$.

Both strategies---structural buffering and behavioural
calibration---are temporary scaffolding: bones harden
and social coupling decreases as the neural network
completes self-calibration.


\subsection{Aging pain immunity}
\label{sec:aging-pain}

If impedance drifts upward with age, why no chronic pain?
For slow drift with healthy C2 ($\eta > 0$),
the impedance-remodeling update continuously re-matches:
$\Gamma^2 \sim (\Delta Z / 2\bar{Z})^2 \lesssim 10^{-6}$.
Pain triggers only in three regimes:
\begin{enumerate}
  \item \textbf{Acute injury} ($dZ/dt \gg \eta$):
    instantaneous discontinuity.
  \item \textbf{C2 failure} ($\eta \to 0$):
    avascular tissues (cartilage) cannot remodel
    $\to$ osteoarthritis (Paper~5~\cite{Paper5}).
  \item \textbf{Late senescence}
    ($\delta D(t) > \eta\Sigma\Gamma^2$):
    Coffin--Manson damage exceeds repair capacity.
\end{enumerate}


% ============================================================
%  SECTION 13 — SCAR TISSUE
% ============================================================
\section{Scar Tissue as Permanent Impedance Debt}
\label{sec:scar-debt}

A wound creates acute localised impedance debt:
$D_{Z,\text{wound}} = \int_0^{T_{\text{repair}}}
\Gamma_{\text{wound}}^2\,P_{\text{in}}\,dt$.
Repair progressively reduces $|\Gamma_{\text{wound}}|$,
but scar tissue never achieves $\Gamma = 0$:
the residual mismatch $\Gamma_{\text{scar}}^2 > 0$
is a \emph{permanent, non-dischargeable impedance debt}
analogous to the irreducible $A_{\text{cut}}$ of Paper~1.

Clinical correlates: scar tenderness (persistent $\Gamma^2$),
contracture (ongoing minimisation), keloid formation
(overshoot to opposite-sign mismatch),
and cumulative surgical debt (raised $\Sigma\Gamma^2$ floor).

\begin{remark}[Scar debt and the bathtub curve]
An individual accumulating many scars shifts baseline
$\Sigma\Gamma^2$ upward, raising the bathtub-curve floor
---a mechanistic link between cumulative trauma and
accelerated aging.
\end{remark}


% ============================================================
%  SECTION 14 — TESTABLE PREDICTIONS
% ============================================================
\section{Testable Predictions}
\label{sec:predictions}

\begin{enumerate}
  \item \textbf{Universality}: Every organism possessing
    neurons must exhibit sleep-like quiescence (sponges excepted).

  \item \textbf{Terrestrial $>$ aquatic sleep complexity}:
    Terrestrial species display more structured sleep
    architecture than aquatic ectotherms of comparable
    neural complexity.

  \item \textbf{Preterm fragility}: Preterm neonates show
    more fragile sleep--thermoregulation coupling.

  \item \textbf{REM as $Z$-map reorganisation}:
    During REM, thermoregulation is suspended for
    unconstrained impedance exploration.

  \item \textbf{Cooling reduces $D_Z$ rate}:
    Lower brain temperature extends tolerable wake duration.

  \item \textbf{Chronic vascular disease as undischarged $D_{Z,v}$}:
    Dual-network cascade predicts that chronic vascular
    debt amplifies neural sleep need.

  \item \textbf{Bathtub curve}: Cognitive performance follows
    the six-phase trajectory without fitted functions.

  \item \textbf{Optimal arousal depends on task difficulty}:
    Yerkes--Dodson optimum shifts with $\Gamma_{\mathrm{env}}$.

  \item \textbf{Depression as high-$\Gamma$, low-$d\Gamma/dt$}:
    Neural network mismatched but unable to learn.

  \item \textbf{PTSD as undamped $\Gamma^2$ spike}:
    Massive mismatch that learning cannot resolve.

  \item \textbf{Dementia as accelerated neural aging}:
    $\delta_n D_n(t)$ overwhelms learning prematurely,
    often preceded by vascular aging cascade.
\end{enumerate}


% ============================================================
%  SECTION 15 — COMPUTATIONAL VERIFICATION
% ============================================================
\section{Computational Verification}
\label{sec:computation}

All simulations use the Alice Gamma-Net framework
with C1--C3 enforced at every time step.

\subsection*{Part~A experiments}

\subsubsection{Experiment~A1: $D_Z$ accumulation and
exponential discharge}

100 wake ticks ($\langle\Gamma^2\rangle \approx 0.05$)
followed by 110 sleep ticks
($\beta_{\mathrm{recal}} = 0.08$).
$D_Z$ rises linearly during wake and decays exponentially
during sleep.
Measured $\beta_{\mathrm{eff}} = 0.078$ matches input
to within 2\%.

\subsubsection{Experiment~A2: Temperature modulation of
$T_{\mathrm{wake}}^{\max}$}

Three simulations with $\tau \in \{0.5, 1.0, 1.5\}$.
Sleep-onset tick ratios: $2.00 : 1.00 : 0.67$---matching
the $1/\tau$ prediction to within 0.2\%.

\subsubsection{Experiment~A3: Infant versus adult profiles}

Infant ($\langle\Gamma^2\rangle = 0.12$, 60/120 wake/sleep)
versus adult ($\langle\Gamma^2\rangle = 0.04$, 120/60).
Infant requires $2.0\times$ more sleep ticks;
the 16\,h/8\,h partition matches neonatal
epidemiology~\cite{Porkka2011}.

\subsubsection{Experiment~A4: Sleep deprivation rebound}

Five deprivation durations (50--250 ticks).
Recovery scales linearly (slope 0.09, $R^2 = 0.93$);
N3 fraction rises from 4.8\% to 17.4\%---consistent with
Process~S dynamics~\cite{Borbely1982}.

\subsubsection{Experiment~A5: $D_{\mathrm{micro}}/D_{\mathrm{thermal}}$
decomposition}

Four environments ($\kappa \in \{1.0, 0.5, 0.2, 0.04\}$).
At $\kappa = 1.0$ (ocean): $D_{\mathrm{thermal}} = 0$.
At $\kappa = 0.04$ (land): $D_{\mathrm{thermal}}/D_{\mathrm{total}} = 49\%$.
Centralisation threshold crossed at $\kappa \approx 0.50$.

\subsubsection{Experiment~A6: Tidal-zone impedance matching}

Three strategies (direct, linear, logarithmic gradient):
$\mathcal{A}[\Gamma] = 0.852, 0.115, 0.053$.
Logarithmic gradient (tidal-zone profile) reduces action
to 6.2\% of direct jump; transmission rises from 14.8\%
to 94.9\%.

\subsubsection{Experiment~A7: Distributed vs.\ centralised
architecture}

$N = 200$ nodes, varied $\kappa$.
Distributed self-cooling suffices at $\kappa \ge 0.50$
(octopus zone); centralised hub wins below
$\kappa_{\mathrm{cross}} \approx 0.30$.
All terrestrial organisms ($\kappa \approx 0.04$) sit
deep in the centralised-required zone.

\subsubsection{Experiment~A8: Dual-field morphogenesis}

1-D PDE simulation:
(a)~Turing patterns in bone ($D_\rho/D_Z = 500$):
$16.4\times$ variance ratio;
(b)~Four-tissue unification: bone, neuron, muscle, liver
all converge under the same PDE with different parameters;
(c)~Ischemia ($I_{\mathrm{blood}} = 0$):
$4.77\times$ $\Gamma^2$ ratio vs.\ healthy tissue.
C1 holds to $< 10^{-15}$ at every grid point.

\subsubsection{Experiment~A9: Embryogenesis verification}

Four scenarios from $Z \approx 0$, $\rho = 0$:
no mother ($\Gamma^2 = 0.937$, failed),
with mother ($\Gamma^2 = 0.037$, developed),
full-term birth ($\Gamma^2 = 0.029$, viable),
premature birth ($\Gamma^2 = 0.038$, permanent deficit).
Confirms material-bottleneck theorem \emph{ab initio}.


\subsection*{Part~B experiments}

\subsubsection{Experiment~B1: Bathtub curve emergence}

Lifecycle Equation integrated for 5000 ticks.
$\Sigma\Gamma^2$ drops from 0.8 to $< 0.2$ by tick 600,
plateaus at 0.1--0.2 for ticks 1000--3000,
rises above 0.3 after tick 4000.
All six developmental phases detected in sequence.

\subsubsection{Experiment~B2: Yerkes--Dodson inverted-U}

Arousal swept from 0 to 1 at fixed $\Gamma_{\mathrm{env}} = 0.5$.
Performance peaks at intermediate arousal ($\sim$0.3--0.5);
degrades at both extremes.

\subsubsection{Experiment~B3: Repeated exposure adaptation}

Fixed stimulus presented 200 times.
$\Gamma$ drops from $\sim$0.7 to $< 0.3$;
learning curve shows exponential decay.

\subsubsection{Experiment~B4: Chronic stress impairs learning}

Relaxed ($\sigma = 0$) versus stressed ($\sigma = 0.7$) agents.
Stressed agent: higher $\Sigma\Gamma^2$, $> 2\times$
aging rate, lower learning efficiency.

\subsubsection{Experiment~B5: Sleep--aging decomposition}

2000-tick simulation with periodic sleep (96 wake / 32 sleep).
Sleep reduces $\Sigma\Gamma^2$ each night;
baseline drifts upward (irreversible plastic accumulation).

\subsubsection{Experiment~B6: Curiosity-driven exploration}

Agent in boring environment.
Boredom accumulates; spontaneous exploration triggers
above threshold; post-exploration $\Gamma^2$ lower than
pre-exploration.

\subsubsection{Experiment~B7: Metacognition System~1/2 switching}

Tasks of varying difficulty.
Easy ($\Gamma_{\mathrm{think}} < 0.3$): System~1;
Hard ($\Gamma_{\mathrm{think}} > 0.6$): System~2;
hysteresis in switching thresholds confirmed.

\subsubsection{Experiment~B8: Novel stimulus shock and recovery}

Well-adapted agent exposed to drastically novel environment.
$\Sigma\Gamma^2$ spikes, recovers exponentially;
post-recovery level slightly lower (net learning).

\subsubsection{Experiment~B9: Full lifecycle simulation}

Complete lifecycle with all three forces, sleep cycles,
stress episodes, and novelty variation.
Six phases emerge; $\Sigma T$ peaks in adulthood;
$T_{\mathrm{steady}}$ stabilises during adult phase.


% ============================================================
%  SECTION 16 — DISCUSSION
% ============================================================
\section{Discussion}
\label{sec:discussion}

\subsection{Methodological caveat}

All 18 experiments use modules built from the $\Gamma$-Net
axioms.
When the simulators reproduce the theory's predictions,
this confirms code--equation consistency, not empirical
truth.
True validation requires longitudinal impedance measurements
(EEG coherence from birth to senescence), comparative
polysomnography across phyla, and prospective clinical tests.
We flag this explicitly so readers assess computational
results as \emph{consistency checks}, not empirical evidence.

\subsection{The brain as exaptation}

The brain evolved as a centralised thermal manager for
impedance-debt dissipation, necessitated by the loss of
the oceanic heat sink.
Cognition, memory, and consciousness are
exaptations~\cite{Torday2016}---secondary functions
exploiting infrastructure built for thermodynamic
housekeeping.
This resolves why the preoptic area controls \emph{both}
sleep and thermoregulation: they were the brain's
original function.

\subsection{Connection to Papers~1--2}

Paper~1 established the \emph{spatial} structure
(relay nodes, fractal topology, dimensional cut cost,
irreducibility).
Paper~2 established the \emph{vascular} extension
(Murray's Law, dual-network organ health,
positive-feedback cascade).
This paper extends both to \emph{time}:
Part~A derives daily dynamics (impedance debt, sleep,
morphogenesis, evolution);
Part~B derives lifecycle dynamics (learning--novelty--aging
competition, emotion, bathtub curve, death).
Together: the static architecture (Paper~1) sets channel
structure $\{Z_i\}$; the vascular model (Paper~2) provides
material coupling $I_{\mathrm{blood}} = (1-\Gamma_v^2)Q_0$;
Part~A determines daily sleep--wake regulation;
Part~B determines the developmental arc of both networks.
All layers derive from a single principle:
$\mathcal{A}[\Gamma] = \int \Sigma|\Gamma|^2\,dt \to \min$.

\subsection{Unification of developmental psychology}

The Lifecycle Equation reduces the entire developmental
trajectory to three physical forces on one state variable.
Piaget's stages are phase transitions of $\Sigma\Gamma^2$;
Vygotsky's ZPD is the $\Gamma_{\mathrm{env}}$ range where
learning can reduce $\Sigma\Gamma^2$ in one session;
critical periods are intervals of enhanced $\eta(t)$.

\subsection{Emotion as physics, not computation}

Dopamine, cortisol, curiosity, and boredom are computed
from $\Gamma^2$ time series in the Alice system.
The cognitive--emotional distinction
dissolves~\cite{Damasio1994}: both are readouts of the
same impedance dynamics at different time scales.

\subsection{Clinical implications}

Depression: high $\Sigma\Gamma_n^2$, low $|d\Gamma_n^2/dt|$
(treatment: increase $\eta_n$ or decrease $\Gamma_{\mathrm{env}}$).
PTSD: undamped $\Gamma^2$ spike (treatment: titrated re-exposure).
ADHD: excess $\gamma$ (stimulants increase effective $\eta$).
Dementia: accelerated $\delta_n D_n$ (often preceded by
vascular cascade).
Hypertension: chronic $\Sigma\Gamma_v^2$ elevation.
Metabolic syndrome: dual-network failure ($\Gamma_n^2$ and
$\Gamma_v^2$ both elevated, $H \to 0$).

\subsection{No-waste corollary}

Every molecular species possesses nonzero $Z$; its presence
modifies $\Gamma$ at every interface.
``Metabolic waste'' is an observational artefact arising
from incomplete knowledge of the full $\Gamma$-network
---consistent with the historical reclassification of
nitric oxide, lactate, bilirubin, uric acid, and adenosine.

\subsection{Dual-field morphogenesis}

The unification of Wolff's law, Davis's law, and
impedance remodeling as parameter instances of a single
PDE (Table~\ref{tab:morphogenesis})
suggests a common thermodynamic origin:
$\mathcal{A}[\Gamma] = \int |\Gamma|^2\,dt \to \min$.
The material bottleneck ($\rho \to 0 \Rightarrow$
only degradation survives) connects directly to
vascular disease via Paper~2:
$\Gamma_v\!\uparrow$ arrests neural plasticity,
structural remodeling, \emph{and} homeostatic regulation
simultaneously.

\begin{table}[h]
  \centering
  \caption{Four-tissue unification: same PDE, different parameters.}
  \label{tab:morphogenesis}
  \begin{tabular}{lcccc}
  \hline
  Tissue & $D_\rho/D_Z$ & $\Gamma^2_i$ & $\Gamma^2_f$ & Reduction \\
  \hline
  Bone (Wolff)     & 500 & 0.0391 & 0.0340 & 13.1\% \\
  Neuron (C2)      &   2 & 0.0391 & 0.0281 & 28.2\% \\
  Muscle (Davis)   &  10 & 0.0391 & 0.0310 & 20.7\% \\
  Liver (hepatic)  & 200 & 0.0391 & 0.0195 & 50.1\% \\
  \hline
  \end{tabular}
\end{table}


% ============================================================
%  SECTION 17 — CONCLUSION
% ============================================================
\section{Conclusion}
\label{sec:conclusion}

We have developed the complete temporal theory of
impedance networks in two complementary parts.

\textbf{Part~A:}
Impedance debt $D_Z = \int |\Gamma|^2 P_{\mathrm{in}}\,dt$
accumulates during wakefulness and is discharged by sleep.
Its decomposition into $D_{\mathrm{micro}}$ and
$D_{\mathrm{thermal}}$ explains:
cnidarian sleep without a brain (ocean removes
$D_{\mathrm{thermal}}$);
convergent brain evolution (loss of heat sink forces
centralisation);
the tidal zone as minimum-action evolutionary path
($\mathcal{A}[\Gamma]$ reduced 93.8\%);
and embryogenesis as thermodynamic causal chain.
A dual-field PDE unifies four tissue-remodeling laws.

\textbf{Part~B:}
The Lifecycle Equation
(Eq.~\eqref{eq:lifecycle})
reduces the entire cognitive trajectory to three forces:
learning, novelty, aging.
Emotion is impedance-mismatch readout at different time scales.
The bathtub curve emerges without fitted functions.
Aging is Coffin--Manson fatigue with Arrhenius acceleration;
atherosclerosis is vascular fatigue;
death is positive-feedback collapse.
Adulthood has a physical definition via the
Equilibrium Equation~\eqref{eq:equilibrium}.

Together with Papers~1--2, this paper completes a four-layer
description of biological intelligence: spatial architecture,
vascular topology, daily maintenance, and the full arc of a life
---all from $\mathcal{A}[\Gamma] \to \min$.
Eighteen computational experiments confirm internal
consistency; empirical falsification via longitudinal
impedance measurements and comparative polysomnography
remains the essential next step.


% ============================================================
%  ACKNOWLEDGMENTS
% ============================================================
\begin{acknowledgments}
The author thanks the open-source scientific Python community
(NumPy, SciPy, pytest) whose tools made the Alice $\Gamma$-Net
simulator possible.
\end{acknowledgments}

\section*{Data and Code Availability}
All simulation code, including the
\texttt{SleepPhysicsEngine},
\texttt{ImpedanceDebtTracker},
\texttt{LifecycleEquationEngine},
\texttt{PinchFatigueEngine},
\texttt{CuriosityDriveEngine},
and \texttt{MetacognitionEngine}
modules, is available at
\url{https://github.com/cyhuang76/alice-gamma-net}
under AGPL-3.0.
Neonatal EEG data sourced from the PhysioNet PICS archive
(open access).


% ============================================================
%  REFERENCES
% ============================================================
\begin{thebibliography}{99}

\bibitem{Paper1}
Hsi-Yu~Huang,
``Minimum Reflection Action on impedance networks:
Irreducibility of dimensional cut cost and fractal scaling,''
Paper~1 of this series, Zenodo (2026).

\bibitem{Paper2}
Hsi-Yu~Huang,
``Dual impedance networks: A unified neural--vascular organ model
from the Minimum Reflection Action,''
Paper~2 of this series, Zenodo (2026).

\bibitem{Paper5}
Hsi-Yu~Huang,
``Clinical validation: Nineteen tissues, one law
--- from NHANES population data to ROC analysis,''
Paper~5 of this series, Zenodo (2026).

\bibitem{Appelbaum2026}
L.~Appelbaum \emph{et al.},
``Sleep is driven by neuronal DNA damage and repair
in jellyfish and sea anemones,''
\emph{Nature Communications} (2026).

\bibitem{Nath2017}
R.~D.~Nath \emph{et al.},
``The jellyfish \textit{Cassiopea} exhibits a sleep-like state,''
\emph{Current Biology} \textbf{27}, 2984 (2017).

\bibitem{Moroz2012}
L.~L.~Moroz,
``Phylogenomics meets neuroscience:
How many times might complex brains have evolved?''
\emph{Acta Biologica Hungarica} \textbf{63}, 3 (2012).

\bibitem{Karbowski2009}
J.~Karbowski,
``Thermodynamic constraints on neural dimensions, firing rates,
brain temperature and size,''
\emph{J.\ Comput.\ Neurosci.} \textbf{27}, 415 (2009).

\bibitem{Porkka2011}
T.~Porkka-Heiskanen and A.~V.~Kalinchuk,
``Adenosine, energy metabolism and sleep homeostasis,''
\emph{Sleep Medicine Reviews} \textbf{15}, 123 (2011).

\bibitem{Borbely1982}
A.~A.~Borb\'{e}ly,
``A two process model of sleep regulation,''
\emph{Human Neurobiology} \textbf{1}, 195--204 (1982).

\bibitem{Wolff1892}
J.~Wolff,
\emph{Das Gesetz der Transformation der Knochen}
(Hirschwald, Berlin, 1892).

\bibitem{Turing1952}
A.~M.~Turing,
``The chemical basis of morphogenesis,''
\emph{Phil.\ Trans.\ R.\ Soc.\ Lond.\ B}
\textbf{237}, 37--72 (1952).

\bibitem{Torday2016}
J.~S.~Torday and W.~B.~Miller,
``On the evolution of the mammalian brain,''
\emph{Frontiers in Systems Neuroscience} \textbf{10}, 31 (2016).

\bibitem{Coffin1954}
L.~F.~Coffin,
``A study of the effects of cyclic thermal stresses
on a ductile metal,''
\emph{Trans.\ ASME}\ \textbf{76}, 931 (1954).

\bibitem{Schultz1997}
W.~Schultz, P.~Dayan, and P.~R.~Montague,
``A neural substrate of prediction and reward,''
\emph{Science}\ \textbf{275}, 1593 (1997).

\bibitem{YerkesDodson1908}
R.~M.~Yerkes and J.~D.~Dodson,
``The relation of strength of stimulus to rapidity
of habit-formation,''
\emph{J.\ Comp.\ Neurol.\ Psychol.}\ \textbf{18}, 459 (1908).

\bibitem{Kahneman2011}
D.~Kahneman,
\emph{Thinking, Fast and Slow}
(Farrar, Straus and Giroux, 2011).

\bibitem{Klutke2003}
G.~A.~Klutke, P.~C.~Kiessler, and M.~A.~Wortman,
``A critical look at the bathtub curve,''
\emph{IEEE Trans.\ Reliab.}\ \textbf{52}, 125 (2003).

\bibitem{Kail1991}
R.~Kail,
``Developmental change in speed of processing during childhood
and adolescence,''
\emph{Psychol.\ Bull.}\ \textbf{109}, 490 (1991).

\bibitem{McEwen1998}
B.~S.~McEwen,
``Stress, adaptation, and disease: Allostasis and allostatic load,''
\emph{Ann.\ N.Y.\ Acad.\ Sci.}\ \textbf{840}, 33 (1998).

\bibitem{Salthouse2010}
T.~A.~Salthouse,
\emph{Major Issues in Cognitive Aging}
(Oxford University Press, 2010).

\bibitem{Huttenlocher1979}
P.~R.~Huttenlocher,
``Synaptic density in human frontal cortex---developmental
changes and effects of aging,''
\emph{Brain Res.}\ \textbf{163}, 195 (1979).

\bibitem{Piaget1952}
J.~Piaget,
\emph{The Origins of Intelligence in Children}
(International Universities Press, 1952).

\bibitem{Vygotsky1978}
L.~S.~Vygotsky,
\emph{Mind in Society: The Development of Higher
Psychological Processes}
(Harvard University Press, 1978).

\bibitem{Hensch2005}
T.~K.~Hensch,
``Critical period plasticity in local cortical circuits,''
\emph{Nat.\ Rev.\ Neurosci.}\ \textbf{6}, 877 (2005).

\bibitem{Damasio1994}
A.~R.~Damasio,
\emph{Descartes' Error: Emotion, Reason, and the Human Brain}
(Avon Books, 1994).

\bibitem{Walker2017}
M.~P.~Walker,
\emph{Why We Sleep: Unlocking the Power of Sleep and Dreams}
(Scribner, 2017).

\bibitem{Ohayon2004}
M.~M.~Ohayon \emph{et al.},
``Meta-analysis of quantitative sleep parameters from
childhood to old age,''
\emph{Sleep}\ \textbf{27}, 1255 (2004).

\bibitem{Mander2017}
B.~A.~Mander, J.~R.~Winer, and M.~P.~Walker,
``Sleep and human aging,''
\emph{Neuron}\ \textbf{94}, 19 (2017).

\bibitem{Epel2004}
E.~S.~Epel \emph{et al.},
``Accelerated telomere shortening in response to life stress,''
\emph{Proc.\ Natl.\ Acad.\ Sci.\ U.S.A.}\ \textbf{101},
17312 (2004).

\bibitem{sprint2019}
SPRINT MIND Investigators,
``Effect of intensive vs standard blood pressure control on
probable dementia,''
\emph{JAMA}\ \textbf{321}, 553 (2019).

\bibitem{Lupien2009}
S.~J.~Lupien \emph{et al.},
``Effects of stress throughout the lifespan on the brain,
behaviour and cognition,''
\emph{Nat.\ Rev.\ Neurosci.}\ \textbf{10}, 434 (2009).

\bibitem{Alter2009}
A.~L.~Alter and D.~M.~Oppenheimer,
``Uniting the tribes of fluency to form a metacognitive nation,''
\emph{Pers.\ Soc.\ Psychol.\ Rev.}\ \textbf{13}, 219 (2009).

\bibitem{Harris2020}
C.~R.~Harris \emph{et al.},
``Array programming with NumPy,''
\emph{Nature}\ \textbf{585}, 357 (2020).

\bibitem{Sorce1985}
J.~F.~Sorce \emph{et al.},
``Maternal emotional signaling: its effect on the
visual cliff behavior of 1-year-olds,''
\emph{Dev.\ Psychol.}\ \textbf{21}, 195 (1985).

\bibitem{Hebb1949}
D.~O.~Hebb,
\emph{The Organization of Behavior}
(Wiley, New York, 1949).

\end{thebibliography}

\end{document}
